% paper3.tex --- Arithmetic Landscape Theory, paper 3.
% SCAFFOLDING ONLY (round 86). Section structure and the verified data are in place;
% the theorem skeleton is fable-5's next deliverable (spec_t3rigid_r085.md).
% Nothing here is a claim-language upgrade: every statement carries its status explicitly.
\documentclass[11pt,a4paper]{article}
\usepackage[T1]{fontenc}
\usepackage{lmodern}
\usepackage{amsmath,amssymb,amsthm}
\usepackage{mathtools}
\usepackage{booktabs}
\usepackage{graphicx}
\usepackage[margin=28mm]{geometry}
\usepackage{microtype}
\usepackage[colorlinks,citecolor=blue,linkcolor=blue,urlcolor=blue]{hyperref}
\usepackage{xcolor}

\theoremstyle{plain}
\newtheorem{theorem}{Theorem}[section]
\newtheorem{proposition}[theorem]{Proposition}
\newtheorem{lemma}[theorem]{Lemma}
\newtheorem{corollary}[theorem]{Corollary}
\theoremstyle{definition}
\newtheorem{definition}[theorem]{Definition}
\newtheorem{problem}[theorem]{Problem}
\newtheorem{conjecture}[theorem]{Conjecture}
\theoremstyle{remark}
\newtheorem{remark}[theorem]{Remark}

\newcommand{\lm}{\mathrm{lm}}
\newcommand{\Lean}[1]{\texttt{#1}}
\newcommand{\TODO}[1]{\textcolor{red}{[\textbf{TODO} #1]}}
\newcommand{\STATUS}[1]{\textcolor{blue}{\textsf{[#1]}}}

\title{The transfer function of subset-sum landscapes:\\
rigidity of the gap series off centre\\[4pt]
\normalsize\textsf{(working title --- scaffolding, round 88)}}
\author{Kentaro Amauchi\thanks{Independent researcher.}}
\date{\today}

\begin{document}
\maketitle

\begin{abstract}
\TODO{written last, after the theorem statements are fixed.} The three things this
paper is to establish: (i) that $\lm_A(n)/\deg_A(n)\to\Gamma(A)$ holds for every fixed
target fraction $\rho$, not only at the centre --- so that $\Gamma$ is an invariant of the
landscape and not of one point of it; (ii) that the ratio is a universal functional of the
\emph{local logarithmic slope} of the representation count alone, through an explicit
transfer function $\Phi$; and (iii) that the hypothesis the first statement needs is exactly
the convergence of the series that gives the second one --- with a profile for which it fails.
\end{abstract}

\section{Introduction}\label{sec:intro}

\TODO{full introduction.} Skeleton of the argument, and what each part rests on:

\begin{itemize}
\item The companion papers establish $\lm/\deg\to\Gamma$ \emph{at the centre}
      $n=\lfloor T/2\rfloor$: paper~1 reduces it to a flatness hypothesis and paper~2
      discharges that hypothesis unconditionally. The obvious question --- whether $\Gamma$ is
      a property of the landscape or of its centre --- is what this paper answers.
\item The combinatorial layer (classification, exact stratification, sandwich) was proved
      \emph{for every target} in paper~1 and is formally verified; nothing there needs redoing.
      \STATUS{Lean-verified}
\item The minor-arc estimates of paper~2 bound $|G_B(\theta)|$ as a function of $\theta$ only.
      They do not refer to the extraction point $m$, which enters
      $r_B(m)=\frac1{2\pi}\int F_B(\theta)e^{-im\theta}d\theta$ only through a phase. They
      therefore transfer to off-centre targets unchanged. \STATUS{quoted from paper 2}
\item What is new is the middle: a tilted local limit theorem and a cancellation lemma for the
      odd orders in $\lambda$. \TODO{fable-5, \texttt{spec\_t3rigid\_r085.md}}
\item One consequence of quoting rather than reproving deserves saying at the outset: \emph{the
      only ineffective constant in the trilogy is one and the same Siegel--Walfisz constant,
      and it enters each paper at exactly one point} --- here, part (b) of
      Lemma~\ref{lem:eta}. See \S\ref{sec:scope}.
\end{itemize}

\section{Notation, and what is inherited}\label{sec:notation}

\TODO{transcribe the definitions from paper 1 in a form that does not assume $n=T/2$.}
Fix throughout an increasing sequence $A=(a_1,\dots,a_k)$ of odd positive integers,
$T=\sigma(A)$, $S_2=\sum a_i^2$, $V=S_2/4$, and a target $n$. Write $\rho=n/T$ and
$x=\tfrac12-\rho$. For $d\ge1$ put
\[
  N_d=\#\{i:a_i\le2d\},\qquad
  \sigma_d=\!\!\sum_{a_i\le 2d}\!\! a_i,\qquad
  s_d=\!\!\sum_{a_i\le 2d}\!\! a_i^{2},\qquad
  \delta_d=d+\tfrac{\sigma_d}{2},
\]
and let $B_d=\{a_i:a_i>2d\}$ be the layer-$d$ ensemble, of variance
$V_d=(S_2-s_d)/4=V-s_d/4$.

\begin{definition}[local logarithmic slope]\label{def:slope}
$\lambda(n)$ is the logarithmic derivative of the representation count at the target,
\[
  \lambda(n)\;=\;-\frac{d}{dm}\log r_A(m)\Big|_{m=n}
  \;=\;\tfrac14\log\frac{r_A(n-2)}{r_A(n+2)}+O(\cdot).
\]
\end{definition}

\begin{remark}[why the slope and not the offset]\label{rem:whyslope}
The Gaussian approximation $\lambda\approx(T/2-n)/V$ is what one writes down first, and it is
wrong at the 10\% level for the purposes of this paper. \TODO{write out: the substitution turns
a $15$--$18\%$ drift in the measured coefficient into $0.03\%$, and it is the reason an earlier
draft of this programme reported a spurious non-integer exponent.} This is not a cosmetic
choice: every statement below is false at the stated precision if $(T/2-n)/V$ is substituted
for $\lambda$.
\end{remark}

\begin{definition}[the tilt]\label{def:tilt}
$P_s$ is the product measure on subsets under which $a_i$ is present with probability
$p_{a}=(1+e^{sa})^{-1}$, and $s=s(n)>0$ is chosen so that $\sum_a a\,p_a=n$; thus $P_s$ is the
exponential tilt of the uniform measure whose mean sits at the target.
\end{definition}

\begin{remark}[the tilt and the slope agree to first order]\label{rem:tiltslope}
$s$ and $\lambda$ are the same quantity to leading order but not exactly: $s$ makes the
tilted \emph{mean} equal $n$, whereas $\lambda$ is the slope of $\log r_A$ \emph{at} $n$, and
the two differ by the mean--mode gap. Measured on three profiles,
\[
   \frac{s(n)}{\lambda(n)}\;=\;1+\frac{c_A}{k}+O(k^{-2}),
   \qquad c_A\approx1.85\ (\text{odds}),\ 2.03\ (\text{primes}),\ 2.87\ (\text{squares}),
\]
stable in $x$ and clean in $k$ over $k=100,\dots,220$.
\STATUS{experimentally confirmed; \texttt{tilt\_r088}}
The statements below are formulated in $\lambda$, and \emph{not} in $s$: substituting $s$
multiplies the residual of Table~\ref{tab:one} by three, exactly as the $\lambda^{2}$ leading
behaviour of $\Phi$ predicts. \STATUS{experimentally confirmed; \texttt{sres\_r088}}
\end{remark}

\section{Rigidity}\label{sec:rigid}

\begin{theorem}[rigidity of the gap series]\label{thm:rigid}
\STATUS{proof skeleton: spec\_t3rigid\_r087 as amended in \S\ref{sec:bridge}}
Let $P=P_k$ be the first $k$ odd primes and let $\rho\in(0,1)$ be fixed. Then
\[
  \frac{\lm_P(\lfloor\rho T\rfloor)}{\deg_P(\lfloor\rho T\rfloor)}
  \;\xrightarrow[k\to\infty]{}\;\Gamma(P),
\]
uniformly for $\rho$ in compact subsets of $(0,1)$. For a general profile $A$ the same holds
under hypothesis~\textup{(H)} of \S\ref{sec:counterexample}.
\end{theorem}

\noindent
$\Gamma(A)=\sum_i a_i2^{-i}$ is therefore an invariant of the landscape and not of its centre.
Theorem~\ref{thm:rigid} is the case $\lambda\to0$ of Theorem~\ref{thm:transfer}: at fixed $\rho$
one has $\lambda(n)\asymp x/N\to0$, and $\Phi$ is continuous at $0$ with $\Phi(0)=\Gamma$, so
\[
  \Phi(\lambda)-\Gamma\;=\;O\Bigl(\lambda^{2}\sum_{d\ge1}2^{1-N_d}\delta_d^{2}\Bigr)
  \;\longrightarrow\;0,
\]
the sum being finite precisely by~(H). \TODO{write out the uniformity in $\rho$; it comes from
$\kappa(\rho)$ being bounded below on compacta, not from any new estimate.}

\STATUS{experimentally confirmed} The evidence the theorem is modelled on. With
$\mathrm{ratio}(\rho,k)=[\lm/\deg]/\Gamma(P_k)$, the spread over $\rho$ contracts monotonically:

\begin{center}
\begin{tabular}{lcccccc}
\toprule
$k$ & 40 & 52 & 64 & 76 & 88 & 100\\
\midrule
$\max_\rho-\min_\rho$ & 0.01051 & 0.00283 & 0.00257 & 0.00162 & 0.00113 & 0.00083\\
$\max_\rho|\mathrm{ratio}-1|$ & 0.00921 & 0.00276 & 0.00244 & 0.00156 & 0.00109 & 0.00081\\
\bottomrule
\end{tabular}
\end{center}

\noindent
over $\rho\in[0.20,0.50]$, by exact integer computation. \TODO{full table; state the
granularity guard $\deg\ge10^4$ and the two correctness checks (agreement with brute force over
all $2^k$ subsets; the reflection identity $\lm(n)=\lm(T-n)$, which is exact).}

\section{The transfer function}\label{sec:transfer}

\begin{theorem}[transfer]\label{thm:transfer}
\STATUS{proof skeleton: spec\_t3rigid\_r087 as amended in \S\ref{sec:bridge}}
Let $A=P_k$ be the first $k$ odd primes, or a general profile satisfying~\textup{(H)}, and let
$n=\lfloor\rho T\rfloor$ with $\rho\in(0,1)$ fixed. Then
\begin{equation}\label{eq:phi}
  \frac{\lm_A(n)}{\deg_A(n)}\;=\;\Phi_k\bigl(\lambda(n)\bigr)\;+\;E_k,
  \qquad
  \Phi_k(\lambda)=1+\sum_{d\ge1}2^{1-N_d}\,
     e^{-\lambda^{2}s_d/8}\,\cosh(\lambda\delta_d),
\end{equation}
with $E_k\to0$ as $k\to\infty$, uniformly for $\rho$ in compact subsets of $(0,1)$.
\end{theorem}

\begin{remark}[what the theorem does and does not claim about $E_k$]\label{rem:whatEk}
The assertion is $E_k\to0$ in absolute terms, nothing sharper. The measured rate
$E_k\asymp c_A/k$ of Remark~\ref{rem:onek} is \emph{not} part of the statement; it is quoted as
\STATUS{experimentally confirmed, $k\le220$, three profiles} and stated as an open problem in
\S\ref{sec:scope}. Distinguishing the two matters here more than usual, because the numerical
agreement is far better than what is proved.
\end{remark}

Three features of \eqref{eq:phi} deserve to be stated before any proof.

\begin{remark}[the ratio sees the target only through one number]
$\Phi$ depends on $A$ alone. The target enters only through $\lambda(n)$. This is a stronger
statement than rigidity: rigidity says the limit does not move, \eqref{eq:phi} says how it
fails to move.
\end{remark}

\begin{remark}[$\lambda=0$ recovers paper 1]\label{rem:lam0}
The layer weight $e^{-\lambda^2s_d/8}$ is $1$ at $\lambda=0$, so
$\Phi(0)=1+\sum_{d\ge1}2^{1-N_d}$, which is the window identity
$W_D(A)=\Gamma(A)+(2D+1)2^{-k}$ of paper~1 --- already formally verified.
Numerically, $\Phi(0)-\Gamma(A)$ is $0$, $-4.7\cdot10^{-14}$ and $-2.7\cdot10^{-15}$ for the
three sequences of Table~\ref{tab:one}.
\end{remark}

\begin{remark}[each layer is a smaller ensemble]\label{rem:weight}
The factor $e^{-\lambda^2 s_d/8}$ is not a fudge: it is what remains when the Gaussian
normalisations of $B_d$ and of $A$ fail to cancel, because $V_d=V-s_d/4$. Expanding to second
order gives the coefficient of $\lambda^2$ as $\delta_d^2-s_d/4$ rather than $\delta_d^2$.
\TODO{write out the two lines; note that the correction is termwise smaller than the term it
corrects, $s_d/4\le d\sigma_d/2<\delta_d^2$, so the sum stays positive.}
\end{remark}

\subsection{Verification}\label{sec:table1}

\STATUS{experimentally confirmed} Table~\ref{tab:one} evaluates \eqref{eq:phi} at the measured
$\lambda(n)$ against the exactly computed $\lm/\deg$.

\begin{table}[htbp]
\centering
\small
\begin{tabular}{lccccccc}
\toprule
$x=\frac12-\rho$ & 0.06 & 0.08 & 0.10 & 0.15 & 0.20 & 0.25 & 0.30\\
\midrule
odds, $k=220$      & 0.84\% & 0.85\% & 0.85\% & 0.86\% & 0.88\% & 0.92\% & 0.97\%\\
squares, $k=170$   & 1.66\% & 1.67\% & 1.68\% & 1.70\% & 1.74\% & 1.80\% & 1.89\%\\
primes, $k=220$    & 0.91\% & 0.92\% & 0.92\% & 0.94\% & 0.96\% & 0.99\% & 1.05\%\\
\bottomrule
\end{tabular}
\caption{Residual $|\,\mathrm{measured}-\Phi(\lambda)\,|$ \textbf{as a fraction of the centred
signal} $\lm/\deg-[\lm/\deg]_{n=T/2}$. Absolute residuals are $10^{-5}$ to $10^{-6}$; quoting
them as significant figures of the ratio would be meaningless, since the signal itself lives in
the fifth and sixth digit.}
\label{tab:one}
\end{table}

\begin{remark}[why the residual is quoted this way]\label{rem:signalfraction}
Against the \emph{uncentred} signal the same residuals read $2\%$ to $300\%$, the large values
occurring precisely where the signal is smallest. Neither number is wrong; the centred one is
the one the theorem is about, because $\Phi(0)=\Gamma$ is an identity and the finite-$k$ offset
$[\lm/\deg]_{n=T/2}-\Gamma$ is a separate quantity of size $10^{-6}$.
\TODO{keep this remark: an earlier draft of the programme reported ``agreement to five or six
significant figures'' for a model that was $12\%$ wrong on the signal.}
\end{remark}

\begin{remark}[the residual obeys a $1/k$ law]\label{rem:onek}
Table~\ref{tab:one} is one column of a $k$-scan. Writing $R_A(k)$ for the residual in the
sense of the caption, the measurement over $k=60,\dots,220$ gives
\[
   k\,R_A(k)\;\longrightarrow\;
   \begin{cases}
     1.85\% & \text{odds } (k=60,100,140,180,220:\ 2.29,2.00,1.91,1.87,1.85),\\
     2.01\% & \text{primes } (k=100,\dots,220:\ 2.38,2.06,2.03,2.01),\\
     2.82\% & \text{squares } (k=130,\dots,190:\ 3.78,2.85,2.83,2.82).
   \end{cases}
\]
Two things follow. First, the residual is $O(1/k)$ and therefore vanishes, which is the form
Theorem~\ref{thm:transfer} asserts; the table's $1\%$ is a finite-$k$ number and not a floor.
Second, the squares row is the loosest of the three because $c_A$ is larger, not because the
decay is different: squares obey the same law with a constant $1.52$ times that of the odd
numbers. The constants $c_A$ coincide numerically with those of
Remark~\ref{rem:tiltslope} --- the same mean--mode gap seen twice --- but substituting $s$ for
$\lambda$ does not remove the residual, so the coincidence is a shared order of magnitude and
not a shared cause. \STATUS{experimentally confirmed; \texttt{tilt\_r088}, \texttt{sres\_r088}}

What they share is a driver. The saddle-point representation gives the exact identity
\[
  s(n)-|\lambda(n)|\;=\;\frac{K'''(s)}{2K''(s)^{2}},
  \qquad K''=\sum a^{2}p_a(1-p_a),\quad K'''=\sum a^{3}p_a(1-p_a)(1-2p_a),
\]
tested pointwise to a relative error $0.54/k$ \STATUS{experimentally confirmed;
\texttt{saddle\_r092}}, and for small tilt $K'''/(2K''^{2})\approx sS_4/S_2^{2}$, so that
$c_A'=k(s/|\lambda|-1)\to kS_4/S_2^{2}$. For $a_i\asymp i^{\alpha}$ this is
$(2\alpha+1)^{2}/(4\alpha+1)$. The measured $c_A$ of this remark agrees with $c_A'$ to within
$1\%$ on four profiles:

\begin{center}
\begin{tabular}{lccccc}
\toprule
 & $\alpha$ & $kS_4/S_2^2$ & $c_A'$ & $c_A$ & $c_A/c_A'$\\
\midrule
odd numbers & 1 & 1.8000 & 1.8525 & 1.8674 & 1.008\\
$a_i\asymp i^{3/2}$ & 3/2 & 2.2794 & 2.3494 & 2.3318 & 0.993\\
squares & 2 & 2.7714 & 2.8615 & 2.8377 & 0.992\\
primes & --- & 1.9731 & 2.0251 & 2.0253 & 1.000\\
\bottomrule
\end{tabular}
\end{center}

\noindent
at $k=220$. \STATUS{experimentally confirmed, four profiles; \texttt{saddle\_r092},
\texttt{alpha\_r092}} \TODO{fable: $c_A=c_A'$ is still a conjecture --- the driver is shared,
the causation is not (substituting $s$ makes the residual worse). State it as a conjecture
with the table as its evidence.}
\end{remark}

\subsection{The bridge to paper 2}\label{sec:bridge}

Both theorems reduce to the same five parts. \textbf{P1} is the exact combinatorial
stratification, proved for every target in paper~1 \STATUS{Lean-verified}; \textbf{P3} is a
local limit theorem for the tilted ensemble of Definition~\ref{def:tilt}; \textbf{P4} is the
assembly, in which the odd orders in $\lambda$ cancel by the reflection symmetry
$\lm(n)=\lm(T-n)$ \STATUS{Lean-verified}; \textbf{P5} is the error budget, which converges
exactly under~(H). Only \textbf{P2} is new: the minor-arc estimates of paper~2 are proved for
the \emph{uniform} measure, and P3 needs them for the tilted one.

Write $\widetilde G_B(\theta)=\prod_{a\in B}\bigl(1-p_a+p_ae(a\theta)\bigr)$ for the tilted
generating function and $S_B(\theta)=\sum_{a\in B}e(a\theta)$ for the linear exponential sum.

\begin{lemma}[the bridge]\label{lem:kappa}
Put $t_a=4p_a(1-p_a)$ and $t_{\min}=\min_{a\in B}t_a$. For every $\theta$,
\[
  \bigl|\widetilde G_B(\theta)\bigr|^{2}
  =\prod_{a\in B}\Bigl(1-t_a\sin^{2}(\pi a\theta)\Bigr)
  \;\le\;\exp\Bigl(-t_{\min}\!\!\sum_{a\in B}\!\sin^{2}(\pi a\theta)\Bigr)
  \;=\;\exp\Bigl(-\tfrac{t_{\min}}{2}\bigl(b-\operatorname{Re}S_B(\theta)\bigr)\Bigr).
\]
Consequently, if $\operatorname{Re}S_B(\theta)\le\eta\,b$ then
$|\widetilde G_B(\theta)|\le\exp\bigl(-t_{\min}(1-\eta)b/4\bigr)$.
\end{lemma}

\begin{proof}
The identity $|1-p+pe(\phi)|^{2}=1-4p(1-p)\sin^{2}(\pi\phi)$ is a computation;
$t_a\ge t_{\min}$ and $1-u\le e^{-u}$ give the middle step; and
$\sum_a\sin^{2}(\pi a\theta)=\tfrac12\bigl(b-\operatorname{Re}S_B(\theta)\bigr)$ is the
double-angle formula summed.
\end{proof}

\begin{remark}[why the bound is additive, and a bound that is not available]\label{rem:noKappa}
It is tempting to transfer multiplicatively, as $|\widetilde G_B(\theta)|\le|G_B(\theta)|^{\kappa}$
with $G_B$ the untilted product, since that would carry every exponent of paper~2 across at
once. \textbf{No such $\kappa>0$ exists.} At any $\theta$ with $a\theta\in\tfrac12+\mathbb Z$ for
some $a\in B$ the untilted factor $\cos(\pi a\theta)$ vanishes, so the right-hand side is $0$,
while the tilted factor equals $\sqrt{1-t_a}>0$. Equivalently
$\log(1-ty)/\log(1-y)\to0$ as $y\to1$. The underlying inequality
$1-ty\le(1-y)^{t}$ is false for every $t\in(0,1)$ --- $s\mapsto(1-y)^{s}$ is convex and
therefore lies \emph{below} its chord $1-sy$, with deficit exactly $-(1-t)$ at $y=1$.
Lemma~\ref{lem:kappa} avoids the issue by never dividing by a quantity that vanishes.
\STATUS{refuted numerically on a $40\,575$-point grid and structurally; \texttt{kappa\_r088}}
\end{remark}

\begin{lemma}[the minor-arc input]\label{lem:eta}
\STATUS{proof skeleton; parts (a) and (c) are quotations, part (b) needs writing}
Let $A=P_k$ and $N=a_k$. Then
\[
  \operatorname{Re}S_A(\theta)\;\le\;\Bigl(\tfrac12+o(1)\Bigr)b
  \qquad\text{uniformly for }\tfrac1{2N}\le\theta\le1-\tfrac1{2N}.
\]
\begin{enumerate}
\item[(a)] \emph{$q\in\{1,2\}$.} Near $\theta\equiv0$ and $\theta\equiv\tfrac12$ partial summation
  against the prime number theorem gives $|S_A(\theta)|\le Cb/(N\|\theta-a/q\|)\le b/2$ once
  $\|\theta-a/q\|\ge1/(2N)$; the excluded interval around $0$ is the principal arc, and around
  $\tfrac12$ one has $\operatorname{Re}S_A\approx-b$, the helpful sign. Elementary; no
  Siegel--Walfisz.
\item[(b)] \emph{$3\le q\le(\log N)^{A_0}$.} Splitting into residue classes and applying
  Siegel--Walfisz, $S_A(a/q+\beta)=\bigl(\mu(q)/\varphi(q)\bigr)W(\beta)+O(\cdot)$ with
  $|W|\le b$, whence $\operatorname{Re}S_A/b\le|\mu(q)|/\varphi(q)\le\tfrac12$. Same toolset as
  paper~2's small-$q$ tier. \TODO{write out; the error term is paper 2's $C_1$, the one
  ineffective constant, and it enters here in the same way and nowhere else.}
\item[(c)] \emph{$q>(\log N)^{A_0}$.} This is the case $n=1$ of Step 2 of
  \textup{prop:deepminor} in paper~2, which gives
  $|S_A(\theta)|\ll N(\log N)^{3-A_0/2}+N^{4/5}(\log N)^{3}$, i.e.\
  $\operatorname{Re}S_A/b\ll(\log N)^{4-A_0/2}=o(1)$ for $A_0>8$.
  \STATUS{quoted from paper 2, verbatim}
\end{enumerate}
\end{lemma}

\begin{remark}[the principal arc, and that the two regimes overlap]\label{rem:principal}
Inside $|\theta|<1/(2N)$ Lemma~\ref{lem:eta} says nothing, and it does not have to: there
$|a\theta|\le\tfrac12$ for every $a\in A$, so $\sin^{2}(\pi a\theta)\ge4a^{2}\theta^{2}$ and
Lemma~\ref{lem:kappa} gives the quadratic bound
$|\widetilde G_A(\theta)|\le\exp(-4t_{\min}V\theta^{2})$ directly. This is already
$\le e^{-Cb}$ for $|\theta|\ge\theta_1:=\sqrt{Cb/(4t_{\min}V)}$, and
\textbf{$\theta_1<1/(2N)$}: with $C=0.01$ the measured $\theta_1N$ is $0.18$--$0.42$ across the
three profiles at $x=0.10$ and $x=0.30$. The two regimes therefore overlap, with the crossover
at $\theta=1/(2N)$ --- which is also, and not by accident, the largest $\theta$ at which the
quadratic bound is valid at all. What remains, $|\theta|\le\theta_1$, is the principal arc,
where P3 supplies the main term rather than a bound.
\STATUS{experimentally confirmed; \texttt{gap\_r090}}
\TODO{\textbf{Seam C.} The bookkeeping of the constant $C$ against P3's local limit theorem is
the one place this could still pinch; it is assigned to the P3 write-out and carried as a named
seam until then. $\theta_1\asymp\sqrt{b/V}$ is $\sqrt b$ times the natural LCLT
width $V^{-1/2}$, which is fine because the integrand is already $e^{-Cb}$ at the edge, but the
error term should be written out rather than asserted.}
\end{remark}

\begin{remark}[what part (c) does \emph{not} need]\label{rem:nocover}
Paper~2 must control every harmonic $n\theta$ with $n\le N_0$, and pays for it with a covering
lemma and the parameter $Q_2=(\log N)^{A_0+A_0'}$. Here only $n=1$ occurs, so the dichotomy in
Lemma~\ref{lem:eta} is on the single denominator $q_1(\theta)$ and no covering lemma is needed.
The quotation in (c) is therefore valid on a strictly larger set of $\theta$ than paper~2's
type~(ii).
\end{remark}

\begin{lemma}[the extremiser of the additive estimate]\label{lem:six}
$\displaystyle\max_{q\ge3}\frac{\mu(q)}{\varphi(q)}=\frac12$, and the maximum is attained at
$q=6$ and at no other $q$.
\end{lemma}

\begin{proof}
$\varphi(q)=2$ holds exactly for $q\in\{3,4,6\}$, where $\mu$ takes the values $-1$, $0$, $+1$;
so among these the maximum is $\mu(6)/\varphi(6)=1/2$. Every other $q\ge3$ has $\varphi(q)\ge4$
and hence $\mu(q)/\varphi(q)\le1/4<1/2$.
\end{proof}

\begin{remark}[the same $q=6$, twice]\label{rem:sixagain}
By Lemma~\ref{lem:eta}(b) the worst point of the additive estimate is therefore $\theta=1/6$,
with $\operatorname{Re}S_A/b\to\tfrac12=\cos60^\circ$. Paper~2's multiplicative estimate has its
worst point at the same $q$: $\sup_qM(q)=M(6)=\sqrt3/2=\cos30^\circ$. \textbf{Same extremiser,
same cause, two functionals} --- both say that the odd primes lie in $\pm1\pmod6$ and nowhere
else, read once through $\prod|\cos|$ and once through $\sum\cos$. Measured on the odd primes
below $10^{7}$, $\operatorname{Re}S_A(1/6)/b=0.499998$ against the predicted $1/2$, and $q=6$
is the maximiser over $2\le q\le16$ at each of $N=10^{5},10^{6},10^{7}$.
\STATUS{experimentally confirmed; \texttt{eta\_r090}}
\end{remark}

\begin{remark}[the honest form of the tilt constant]\label{rem:kapparho}
$t_{\min}=t_{\min}(\rho)$ is positive for every fixed $\rho\in(0,1)$ and bounded below on
compact subsets, which is all Lemma~\ref{lem:kappa} needs. It is \emph{not} governed by the
approximation $sN\approx6x$: that relation is the small-$x$ asymptote only, and the measured
ratio $sN/6x$ is $1.02$--$1.13$ at $x=0.10$ but $1.22$--$1.39$ at $x=0.30$. Measured values of
$t_{\min}$: $0.91$, $0.89$, $0.91$ at $x=0.10$ and $0.36$, $0.28$, $0.34$ at $x=0.30$ for the
odd numbers, squares and primes.
\STATUS{experimentally confirmed; \texttt{tilt\_r088}, \texttt{kappa\_r088}}
\end{remark}

\noindent
With Lemmas~\ref{lem:kappa} and~\ref{lem:eta} the tilted generating function decays like
$\exp(-t_{\min}(1-\eta)b/4)$ off the principal arc --- exponentially in $b$, which is what P5
quotes. The tilted analogue of the $M(q)$ spectrum is lost, and is not needed: the rate enters
Theorem~\ref{thm:transfer} only through $E_k=o(1)$.

\section{Universality of the rate function}\label{sec:universality}

\STATUS{experimentally confirmed, $k=16..56$, three reference sequences}
\TODO{full table.} With $\hat I_A(\rho)=-k^{-1}\log(\deg_A(\lfloor\rho T\rfloor)/2^k)$ and
$D_k=\hat I_P-\hat I_{\mathrm{ref}}$, the difference tends to $0$ for each of three reference
sequences: the odd numbers, the arithmetic progression $6n+5$, and a random odd sequence of the
same range. The sign is a property of the reference, not a universal fact --- it is positive
against the odd numbers, negative against the random sequence, and changes sign near $k=48$
against the progression.

\begin{remark}[the primes sit near the random sequence]\label{rem:random}
At $k=56$ the distance to a random odd sequence is $0.0016$ against $0.0146$ to the arithmetic
progression --- an order of magnitude. \TODO{develop: this is the direct evidence for the
reading that the main term is universal and the arithmetic of the sequence appears only in the
correction.}
\end{remark}

\section{A profile for which the hypothesis fails}\label{sec:counterexample}

The hypothesis Theorem~\ref{thm:rigid} needs is the convergence, with room, of the series in
\eqref{eq:phi} --- qualitatively, that $N_d\to\infty$ fast enough. \TODO{state it precisely
once fable-5's skeleton fixes the form.}

\STATUS{experimentally confirmed} It is not vacuous. For the odd-ified cubes
$a_i=2\lfloor i^3/2\rfloor+1$:

\begin{center}
\begin{tabular}{lcccc}
\toprule
profile & odds & primes & squares & cubes\\
\midrule
$N_d$ at $d=10$ (any $k$) & 10 & 7 & 4 & \textbf{2}\\
$Q(0)=\Gamma^{-1}\sum2^{-N_d}(\delta_d^2-s_d/4)$ & 20.3 & 50.4 & 916 & \textbf{381\,904}\\
\bottomrule
\end{tabular}
\end{center}

\noindent
Only $\{1,9\}$ ever lie below $20$, at any $k$, so $2^{-N_d}$ never damps the growing
$\delta_d^2$ and the series turns over only near $d\approx10^4$. At $k=70$ the finite-size
offset is still $1.75\cdot10^{-2}$ while the signal is $1.7\cdot10^{-11}$: nine orders of
magnitude, so the profile is not merely slow but undecidable at any accessible size.
\TODO{present this as the counterexample that makes the hypothesis non-vacuous, not as a
failure.}

\subsection{Two different kinds of failure}\label{sec:twofailures}

The cubes fail \emph{effectively} and not asymptotically, and the distinction is worth
keeping. Computed directly, $\sum_d2^{-N_d}\delta_d^2$ for the odd-ified cubes settles at
$1.112\cdot10^{7}$ and does not grow with $k$ (values at $k=60,100,140,180$:
$1.112\cdot10^{7}$ from $k\approx140$ on, against $69.5$ for the odd numbers, $266.1$ for
$\alpha=3/2$ and $6473$ for the squares --- each of them $k$-independent as well).
\STATUS{experimentally confirmed; \texttt{alpha\_r092}} So \textup{(H)} is not what the cubes
violate; what they violate is any hope of seeing the asymptotics, the constant being four
orders of magnitude larger than the odd numbers'.

A profile that fails \textup{(H)} outright is available, and on the other side:

\begin{proposition}[no profile below $\alpha=1$]\label{prop:alphalb}
Let $a_i\asymp c\,i^{\alpha}$ be distinct odd integers with $\alpha<1$. Then
$c\gg k^{1-\alpha}$, hence $a_1\to\infty$, hence $N_d=0$ for all $d<a_1/2$ and
\[
  \sum_{d\ge1}2^{-N_d}\delta_d^{2}\;\ge\;\sum_{d<a_1/2}2d^{2}\;\asymp\;a_1^{3}
  \;\longrightarrow\;\infty .
\]
So \textup{(H)} fails for every $\alpha<1$.
\end{proposition}

\begin{proof}
The $k$ values $a_1<\dots<a_k$ are distinct odd integers, so $a_k\ge2k-1$; with
$a_k\asymp ck^{\alpha}$ this forces $c\gg k^{1-\alpha}$. For $d<a_1/2$ no $a_i$ is $\le2d$, so
$N_d=0$, $\sigma_d=0$ and $\delta_d=d$, and the displayed terms are $2d^2$.
\end{proof}

\noindent
Measured on an $\alpha=1/2$ profile, $\sum2^{-N_d}\delta_d^2=4923,\,10245,\,16608,\,24464$ at
$k=60,100,140,180$ --- growing like $k^{3/2}$, as the proposition predicts.
\STATUS{experimentally confirmed; \texttt{alpha\_r092}} The hypothesis therefore cuts the
power profiles at $\alpha=1$ exactly, and the cubes sit inside the admissible range but
beyond the reach of computation.

\section{Honest scope}\label{sec:scope}

\TODO{in the style of paper 1 \S10 and paper 2 \S10.} The classification to fill in:

\begin{itemize}
\item \emph{Formally verified}: the combinatorial layer, inherited from paper~1, for every
      target.
\item \emph{Proved}: \TODO{after fable-5's skeleton.} The minor-arc input is quoted from
      paper~2 and is $m$-independent; paper~2's own ineffectivity through the Siegel--Walfisz
      constant is inherited.
\item \emph{Experimentally confirmed, with the range stated}: rigidity
      (\S\ref{sec:rigid}), the transfer function to $0.8$--$1.9\%$ of the signal
      (Table~\ref{tab:one}), universality (\S\ref{sec:universality}), the profile-dependent
      constant $4(2\alpha+1)/(\alpha+1)$ verified on four profiles.
\item \emph{Not claimed}: the residual $0.8$--$1.9\%$ of Table~\ref{tab:one} is consistent
      with finite size and shrinks in $k$, but we do not claim it is zero.
\end{itemize}

\begin{remark}[the single ineffective constant]\label{rem:oneconstant}
Precisely: every estimate used here is effective except part~(b) of Lemma~\ref{lem:eta}, whose
error term is the Siegel--Walfisz error, and paper~2's appendix identifies the same constant
$C_1$ as the only ineffective one in that paper, entering there at one point as well. Paper~1
is unconditional and effective throughout. \textbf{Ineffective is not conditional}: no
statement here depends on an unproved hypothesis; what is not available is a numerical
threshold beyond which the asymptotics take hold.
\end{remark}

\begin{remark}[approaches that did not work]\label{rem:graveyard}
\TODO{as in paper 2, keep the dead ends with their reasons.} At least: the tilted-measure
($\Gamma_x$) reading of off-centre targets, which the data rejected --- rewriting the collapse
in the exact tilt parameter makes it worse, not better; and the reading of the correction as a
third-cumulant effect, which was an artefact of using the Gaussian slope.
\end{remark}

\section*{Data availability}
\TODO{scripts and logs, as in papers 1 and 2.}

\begin{thebibliography}{9}
\bibitem{P1} K.~Amauchi, \emph{The gap series of an integer sequence: an arithmetic invariant
governing subset-sum landscapes}. \TODO{arXiv number}
\bibitem{P2} K.~Amauchi, \emph{Asymptotic flatness of subset-sum landscapes of primes: the
sub-peak spectrum and the constant $\sqrt3/2$}. \TODO{arXiv number}
\end{thebibliography}

\end{document}
